% Synthetic wide/tall PDFs are prepared by tools/check_theme.py.
\documentclass[language=english,cjk=false]{campusbeamer}
\title{Page styles}
\author{Template check}
\begin{document}
\section{Alpha}
\subsection{Context}
\begin{frame}{Ordinary before}Body before.\end{frame}
\begin{frame}[campusfocus]{Plain focus}Body without a background.\end{frame}
\begin{frame}[campusbackground=build/style-wide.pdf,campusbackgroundopacity=1]{Wide background}
  \uncover<1->{First stage.}\par
  \uncover<2->{Second stage.}
  \note<1>{First background stage.}\note<2>{Second background stage.}
\end{frame}
\begin{frame}[campusbackgroundopacity=.5,campusbackground=build/style-tall.pdf]{Tall background}
  Portrait image, centered and cropped.
\end{frame}
\begin{frame}[campusbackground=build/style-wide.pdf,campusbackgroundopacity=0]{Invisible background}
  Fully transparent background.
\end{frame}
\begin{frame}[fragile,campusfocus]{Fragile focus}
\begin{verbatim}
print("frame keys work with fragile")
\end{verbatim}
\end{frame}
\begin{frame}[campusfocus,allowframebreaks]{Focus continuation}
First page.\par\framebreak
Second page.
\end{frame}
\begin{frame}{Ordinary after}Body after.\end{frame}
\end{document}
